% ---
% id: foundations/lem:least-eigenvalue-bound
% kind: lem
% title: "Least-eigenvalue lower bound from an approximate eigenbasis"
% summary: >-
%   For a Hermitian matrix and any real V whose Gram matrix lies within θ < 1
%   of the identity in Frobenius norm, the Gershgorin bound γ of the real
%   embedding conjugated by V gives λ_min ≥ γ/(1+θ), or γ/(1−θ) if γ < 0. On
%   rational data every step is exact, and V decides tightness only.
% prominence: standard
% status: implemented
% uses: []
% mentions: [def:exact-island]
% assumes: []
% refs: []
% provenance: newly-derived
% implements: certiqs_key.foundations.dyadic:least_eigenvalue_bound
% lean:
% ---

An eigensolver in binary64 returns approximate eigenvectors, and nothing is
established of them by the arithmetic that produced them. The lemma turns any
such basis into a lower bound on the least eigenvalue that holds exactly: it
asks of the basis only a bound on its defect from orthogonality, so the
quality of the basis decides how close the bound lies to the least eigenvalue
and never whether it is one. The real embedding lets the computation run over
the reals, a Hermitian matrix having the Rayleigh quotients of its embedding,
and on rational data every quantity is rational, so the bound is a
certification of the \tref{exact island} (Definition~\ref{def:exact-island}).

\begin{lemma}\label{lem:least-eigenvalue-bound}
Let $n\ge1$, let $R$ be a Hermitian complex $n\times n$ matrix, and let
\begin{equation}\label{eq:real-embedding}
  \widetilde R\;\coloneqq\;
  \begin{pmatrix}\mathrm{Re}\,R&-\mathrm{Im}\,R\\
  \mathrm{Im}\,R&\mathrm{Re}\,R\end{pmatrix}\in\mathbb{R}^{2n\times2n},
\end{equation}
the real and imaginary parts taken entrywise. Let $V$ be any real
$2n\times2n$ matrix, put
\begin{equation*}
  B\coloneqq V^{\transpose}\widetilde R\,V,\qquad
  E\coloneqq V^{\transpose}V-\idon{2n},\qquad
  \gamma\coloneqq\min_{1\le i\le2n}\Big(B_{ii}-\sum_{l\ne i}|B_{il}|\Big),
\end{equation*}
and let $\theta$ be a real number with $0\le\theta<1$ and
$\sum_{i,l}E_{il}^{2}\le\theta^{2}$.
\begin{enumerate}[label=(\roman*),leftmargin=2.4em]
\item $\widetilde R$ is real symmetric, and $\lmin(\widetilde R)=\lmin(R)$.
\item $V$ is invertible, and
\begin{equation}\label{eq:least-eigenvalue-bound}
  \lmin(R)\;\ge\;
  \begin{cases}
    \gamma/(1+\theta), & \gamma\ge0,\\
    \gamma/(1-\theta), & \gamma<0.
  \end{cases}
\end{equation}
\item If the entries of $R$ lie in $\mathbb{Q}(i)$, the entries of $V$ are
  rational and $\theta$ is rational, then $\widetilde R$, $B$ and $E$ have
  rational entries, $\gamma$ is rational, the hypothesis on $\theta$ is a
  comparison of rationals, and the right-hand side of
  \eqref{eq:least-eigenvalue-bound} is a rational number.
\item Applied to $-R$, (ii) bounds the largest eigenvalue of $R$, which is
  $-\lmin(-R)$, from above.
\end{enumerate}
\end{lemma}

\begin{proof}
\emph{(i).} Write $A\coloneqq\mathrm{Re}\,R$ and $C\coloneqq\mathrm{Im}\,R$.
Since $R$ is Hermitian, $\overline{R_{li}}=R_{il}$ for all $i,l$, so
$A_{li}=A_{il}$ and $C_{li}=-C_{il}$: $A$ is symmetric and $C$ is
antisymmetric. Hence
$\widetilde R^{\transpose}=\begin{pmatrix}A^{\transpose}&C^{\transpose}\\
-C^{\transpose}&A^{\transpose}\end{pmatrix}
=\begin{pmatrix}A&-C\\C&A\end{pmatrix}=\widetilde R$. For
$x=\begin{pmatrix}a\\b\end{pmatrix}\in\mathbb{R}^{2n}$, with
$a,b\in\mathbb{R}^{n}$, put $u\coloneqq a+ib\in\mathbb{C}^{n}$. Then
$Ru=(Aa-Cb)+i(Ca+Ab)$ and
\begin{equation*}
  u^{\dagger}Ru
  =a^{\transpose}(Aa-Cb)+b^{\transpose}(Ca+Ab)
  +i\big(a^{\transpose}(Ca+Ab)-b^{\transpose}(Aa-Cb)\big).
\end{equation*}
The number $u^{\dagger}Ru$ is real, $R$ being Hermitian, so it equals its
real part
$a^{\transpose}Aa-a^{\transpose}Cb+b^{\transpose}Ca+b^{\transpose}Ab$, which
is $x^{\transpose}\widetilde Rx$ written out by blocks; and
$u^{\dagger}u=a^{\transpose}a+b^{\transpose}b=x^{\transpose}x$. The map
$x\mapsto u$ is a bijection of $\mathbb{R}^{2n}$ onto $\mathbb{C}^{n}$ that
sends $0$ to $0$, so the set of Rayleigh quotients
$u^{\dagger}Ru/u^{\dagger}u$, $u\ne0$, equals the set of
$x^{\transpose}\widetilde Rx/x^{\transpose}x$, $x\ne0$. By the spectral
theorem the least eigenvalue of a Hermitian matrix is the minimum of its
Rayleigh quotient over non-zero complex vectors, and the least eigenvalue of
a real symmetric matrix, which has a real orthonormal eigenbasis, is the
minimum of its Rayleigh quotient over non-zero real vectors. So
$\lmin(\widetilde R)=\lmin(R)$.

\emph{(ii), the Gram matrix.} For $y\in\mathbb{R}^{2n}$,
$\|Vy\|^{2}=y^{\transpose}V^{\transpose}Vy=\|y\|^{2}+y^{\transpose}Ey$.
By the Cauchy--Schwarz inequality applied to each row of $E$,
$\|Ey\|^{2}=\sum_i\big(\sum_lE_{il}y_l\big)^{2}
\le\sum_i\sum_lE_{il}^{2}\,\|y\|^{2}\le\theta^{2}\|y\|^{2}$, and again by
the Cauchy--Schwarz inequality $|y^{\transpose}Ey|\le\|y\|\,\|Ey\|
\le\theta\|y\|^{2}$. Therefore
\begin{equation}\label{eq:gram-sandwich}
  (1-\theta)\|y\|^{2}\;\le\;\|Vy\|^{2}\;\le\;(1+\theta)\|y\|^{2}
  \qquad\text{for every }y\in\mathbb{R}^{2n}.
\end{equation}
As $\theta<1$, $Vy=0$ forces $y=0$, so the square matrix $V$ is injective
and hence invertible.

\emph{(ii), Gershgorin's bound on the quadratic form.} $B$ is real symmetric,
$B^{\transpose}=V^{\transpose}\widetilde R^{\transpose}V=B$ by (i). For
$y\in\mathbb{R}^{2n}$, using $|y_iy_l|\le\tfrac12(y_i^{2}+y_l^{2})$, so that
$B_{il}y_iy_l\ge-|B_{il}|\,|y_iy_l|\ge-\tfrac12|B_{il}|(y_i^{2}+y_l^{2})$
whatever the signs of $B_{il}$, $y_i$ and $y_l$, and
$|B_{il}|=|B_{li}|$,
\begin{equation*}
  y^{\transpose}By=\sum_iB_{ii}y_i^{2}+\sum_{i\ne l}B_{il}y_iy_l
  \ge\sum_iB_{ii}y_i^{2}-\tfrac12\sum_{i\ne l}|B_{il}|\big(y_i^{2}+y_l^{2}\big)
  =\sum_iy_i^{2}\Big(B_{ii}-\sum_{l\ne i}|B_{il}|\Big)
  \ge\gamma\,\|y\|^{2},
\end{equation*}
the equality because the two halves of the double sum are equal after the
indices are exchanged, $|B_{il}|$ being symmetric in $i$ and $l$.

\emph{(ii), the bound.} Let $x\in\mathbb{R}^{2n}$, $x\ne0$, and put
$y\coloneqq V^{-1}x\ne0$. Then
$x^{\transpose}\widetilde Rx=y^{\transpose}V^{\transpose}\widetilde RVy
=y^{\transpose}By\ge\gamma\|y\|^{2}$, and $\|x\|^{2}=\|Vy\|^{2}$ lies
between $(1-\theta)\|y\|^{2}$ and $(1+\theta)\|y\|^{2}$ by
\eqref{eq:gram-sandwich}. If $\gamma\ge0$, then
$\|y\|^{2}\ge\|x\|^{2}/(1+\theta)$ gives
$\gamma\|y\|^{2}\ge\gamma\|x\|^{2}/(1+\theta)$. If $\gamma<0$, then
$\|y\|^{2}\le\|x\|^{2}/(1-\theta)$, the denominator being positive, gives
$\gamma\|y\|^{2}\ge\gamma\|x\|^{2}/(1-\theta)$. In either case the Rayleigh
quotient $x^{\transpose}\widetilde Rx/x^{\transpose}x$ is at least the
right-hand side of \eqref{eq:least-eigenvalue-bound}. Since $x$ was
arbitrary, so is $\lmin(\widetilde R)$, the minimum of that quotient, and
$\lmin(\widetilde R)=\lmin(R)$ by (i).

\emph{(iii).} The real and imaginary parts of an element of
$\mathbb{Q}(i)$ are rational, so $\widetilde R$ is rational; $B$ and $E$
are sums of products of rationals; $|B_{il}|$, the sums and the minimum
defining $\gamma$ are rational; $\sum_{i,l}E_{il}^{2}\le\theta^{2}$,
$0\le\theta$ and $\theta<1$ compare rationals; and $1\pm\theta>0$ is
rational, so the quotient is.

\emph{(iv).} $-R$ is Hermitian, its real embedding is $-\widetilde R$, and
its eigenvalues are those of $R$ with the sign reversed, so the largest
eigenvalue of $R$ is $-\lmin(-R)$, and a lower bound on $\lmin(-R)$ from
(ii) gives an upper bound on it.
\end{proof}
